\section{Typed Flexible-Arity Slotted E-Graphs}
\label{sec:formalism}

We fix a typed source signature and a choice of source equational laws. The
representation separates three responsibilities: typed maps describe how an open
expression is placed into a caller, ports describe the algebraic structure of
an operator's children, and graph records track equations already proved.

\paragraph{Two source equalities.}
If both operands are unary relations, domain and range restriction select
exactly the atoms in their intersection. This does not replace restriction
of a binary relation: that operation selects tuples by one column.
For a fixed domain $D$, $\forall x\in D.\,\Box P(x)$ and
$\Box(\forall x\in D.\,P(x))$ both require $P(x)$ for every atom in $D$
at every future instant. An empty $D$ is allowed. If $D$ can grow or
shrink, the two formulas inspect different bindings; the rewrite is then
not justified. The appendix gives the Lean statements and negative controls.

\subsection{Contexts, invocations, and ports}

A \emph{typed context} $\Gamma$ is a finite collection of distinct slots, each
carrying a type. A \emph{typed embedding} is an injective, type-preserving map;
a \emph{typed renaming} is an embedding with a typed inverse:
\[
 \TEmb(S,T)=\{m:S\hookrightarrow T\mid
                    \operatorname{ty}(m(x))=\operatorname{ty}(x)\},
\]
and $\TRen(S,T)$ denotes the typed renamings. Identities and composites are
embeddings, and renamings are closed under inverse and composition. Fixing a
linear order on types defines a canonical context $\operatorname{Can}(S)$: list
types in order and index the slots of each type from $0$. This preserves the
slot count of each type without retaining original variable names.

An e-class $a$ has interface $S_a$ and output type $\tau_a$. An
\emph{invocation} of $a$ in caller $\Gamma$ is $m*a$, where
$m\in\TEmb(S_a,\Gamma)$; its type is $\tau_a$. We abbreviate
$\operatorname{id}_{S_a}*a$ by $a$. The map is part of the occurrence: two
invocations of the same class can supply different caller variables.

Each type-instantiated operator $f_\theta$ declares input \emph{ports} using
containers $\Seq$, $\Bag$, or $\Set$ (sequences, multisets, or sets) paired
with a nonempty set $J\subseteq\mathbb N$ of legal arities ($K^+$ for positive
arities, $K^{=k}$ for arity $k$):
\[
 \kappa ::= \texttt{One}(\tau)\mid K^J(\kappa)
 \mid\texttt{Bind}(\tau,\kappa)\mid\texttt{BindBlock}(\beta,\kappa),
 \qquad K\in\{\Seq,\Bag,\Set\}.
\]
Checked syntax retains occurrence vectors even for bag and set ports so
provenance can track duplicate removal and reordering; normalization then
computes the sequence, multiset, or set key. A typed node has the form
\[
 n=f_\theta(q_1,\ldots,q_r),\qquad
 q_i\in\operatorname{Port}_\Gamma(\kappa_i),\qquad
 \operatorname{out}(n)=\tau_{f_\theta}.
\]
Port atoms are free slots, typed de Bruijn bound indices, or invocations.
Binders extend the local de Bruijn scope. A binder-block descriptor $\beta$
records types, domains, multiplicities, disjointness, and dependencies,
together with an admitted subgroup $\operatorname{Aut}(\beta)$ of typed bound
permutations whose elements preserve meaning under the descriptor's conditions.

\paragraph{Licensed container laws.}
A $\Seq$ port requires no sibling equation, $\Bag$ requires commutativity, and
$\Set$ requires commutativity and idempotence. Recursive flattening is licensed
separately: it applies only to a sole root container with matching element and
result types, requires associativity for that exact operator instance, and
requires arity closure:
\begin{align*}
 k\in J,\ 1\le j\le k &\implies j\in J &&(\text{for }\Set),\\
 k,l\in J,\ k>0 &\implies k+l-1\in J &&(\text{for flat}).
\end{align*}
If a flat port admits zero operands, the source theory must also supply a unit.
These laws must be justified by the source theory for the exact operator
instance: conjunction uses a flat $\Set$, equality uses a commutative binary
port without flattening, and integer addition uses a $\Bag$ whose flattening
depends on whether arithmetic is modular or overflow-checked.

\paragraph{Renaming and support.}
The free support $\operatorname{slots}(v)$ collects the free slots in an
expression: a free slot contributes itself, $m*a$ contributes $\operatorname{im}(m)$,
bound indices contribute nothing, and composites take the union of their
children. Any embedding $e\in\TEmb(\Gamma,\Gamma')$ acts componentwise, fixes
bound indices, and maps $m*a$ to $(e\circ m)*a$. By structural induction,
this action preserves schemas and output types and satisfies
\begin{equation}
 \operatorname{slots}(e\cdot v)=e[\operatorname{slots}(v)].
 \label{eq:support-action}
\end{equation}

\subsection{A worked example: sharing without losing variables}
\label{sec:worked-example}

Let class $a$ represent the Alloy relation test $x\in y.\mathsf{follows}$ over
$S_a=\{s:\mathsf{User},\,t:\mathsf{User}\}$, and let the caller have
$x,y:\mathsf{User}$. The two maps
$m(s)=x, m(t)=y$ and $m'(s)=y, m'(t)=x$ give invocations
$m*a=(x\in y.\mathsf{follows})$ and $m'*a=(y\in x.\mathsf{follows})$. Inside
a conjunction $\Set$ port $\{m*a,m'*a\}$, both invocations share class $a$ but
remain distinct elements because $m\neq m'$. If the graph establishes
$s\in t.\mathsf{follows}\iff t\in s.\mathsf{follows}$ and records the swap
$\pi(s)=t, \pi(t)=s$ as a symmetry of $a$, then $m'=m\circ\pi$ puts both
invocations in one local orbit. Normalizing that orbit before forming the set
collapses $\{m*a,m'*a\}$ to one element; under a $\Bag$, both multiplicities
remain, giving multiplicity~2; under a $\Seq$, both positions remain.

A caller renaming acts globally on an entire kernel—renaming $(x,y)$ to
$(u,v)$ preserves the relationship between $m$ and $m'$—whereas class
symmetries act per invocation and binder-block permutations act only within a
quantifier or comprehension block whose descriptor conditions admit them.

\subsection{Graph structure and effective support}

A finite graph stores class interfaces, typed node cases, parent links, and
recorded interface symmetries. A parent link $a\xrightarrow{p}b$ carries
$p\in\TEmb(S_b,S_a)$ between classes of equal output type. The parent forest has
decreasing depths, so leader lookup terminates at a root $\ell$ with composite
embedding $p_a\in\TEmb(S_\ell,S_a)$, sending $m*a$ to $\widehat m*\ell$ where
$\widehat m=m\circ p_a$. Because $p_a$ need not be surjective,
\begin{equation}
 \operatorname{im}(m\circ p_a)\subseteq\operatorname{im}(m),
 \label{eq:find-support}
\end{equation}
so leader normalization can contract the exposed support.

Each leader $\ell$ records a group $G_\ell$ of justified typed symmetries. Two
leader invocations align under a caller renaming $\rho\in\TRen(\Gamma,\Gamma')$
when their leaders coincide and there exists $\pi\in G_\ell$ such that
\begin{equation}
 \rho\circ\widehat m=\widehat m'\circ\pi.
 \label{eq:invocation-alignment}
\end{equation}
The structural equivalence relation compares sequences pointwise, bags by an
occurrence bijection, sets by mutual forward and backward matching, unary
binders by preserving the bound index, and binder blocks via admitted
descriptor permutations, closed under reflexivity, symmetry, and transitivity.
By composing~\eqref{eq:invocation-alignment} with inverses, caller renamings
transport alignments via conjugation and leader symmetries preserve image
supports.

\subsection{Canonical shapes}
\label{sec:algorithm}

Canonicalization first follows parent links and applies licensed local
simplifications to extract a \emph{kernel} $K_n$ over its effective support
$\Delta_n\subseteq\Gamma_n$, retaining the support inclusion
$\iota_n\in\TEmb(\Delta_n,\Gamma_n)$ and finite provenance between $n$ and
$\iota_n\cdot K_n$. Within $K_n$, child and binder orbits are normalized before
aggregating bags or deduplicating sets, and a least representative is chosen
under free-slot renaming.

To state exactness abstractly, a \emph{finite quotient presentation} fixes a
linearly ordered carrier $X$, a local structural generator whose generated
equivalence $\sim$ is decidable, a finite list enumerating every candidate in
$X$ (\emph{carrier completeness}), and invariants asserting that each generator
step preserves support, schema, leaf types, and output type. Define
$Q(x)=\min\{y\in X\mid y\sim x\}$. An \emph{effective-support canonicalizer}
additionally maps each node $n$ to an effective kernel $K_n$ over effective
context $\Delta_n$, an inclusion $\iota_n\in\TEmb(\Delta_n,\Gamma_n)$, a
constructive witness selector for $\sim$, and kernel preservation of support,
schema, leaf types, and output type, such that $K_n$ and $K_{n'}$ are related
by the presentation equivalence if and only if a typed renaming
$\rho\in\TRen(\Delta_n,\Delta_{n'})$ carries a structural derivation
$K_n\Rightarrow_\rho K_{n'}$. Put $\operatorname{Shape}(n)=Q(K_n)$.

\begin{theorem}[Effective-support canonical-shape exactness]
\label{thm:shape-exactness}
Under a finite quotient presentation, $Q(x)\sim x$, $Q(Q(x))=Q(x)$,
$Q(x)=Q(y)\iff x\sim y$, and $Q$ preserves support, schema, leaf types, and
output type. Consequently, for any effective-support canonicalizer as above and
nodes $n,n'$,
\[
 \operatorname{Shape}(n)=\operatorname{Shape}(n')
 \quad\Longleftrightarrow\quad
 \exists\rho\in\TRen(\Delta_n,\Delta_{n'}).\ K_n\Rightarrow_\rho K_{n'},
\]
and either side implies $\operatorname{out}(n)=\operatorname{out}(n')$.
Moreover, the constructed canonical record supplies a witness
$\sigma_n\in\TRen(C_n,\Delta_n)$ (where $C_n=\operatorname{Can}(\Delta_n)$)
and ambient transport $\omega_n=\iota_n\circ\sigma_n\in\TEmb(C_n,\Gamma_n)$.
\end{theorem}
\begin{proof}
For every $x\in X$, the finite orbit $\{y\in X\mid y\sim x\}$ is nonempty because
$x\sim x$, so its minimum $Q(x)$ exists and satisfies $Q(x)\sim x$. If $x\sim y$,
symmetry and transitivity identify their orbits, giving $Q(x)=Q(y)$;
conversely, $Q(x)=Q(y)$ implies $x\sim Q(x)=Q(y)\sim y$. Setting $y=Q(x)$ yields
idempotence. Induction on the reflexive--symmetric--transitive closure lifts
support, schema, leaf-type, and output-type invariance from each generator to
$\sim$, and hence to $Q(x)\sim x$.

Applying this exactness to $K_n$ and $K_{n'}$ and composing with the
presentation-to-derivation equivalence yields
\[
 \operatorname{Shape}(n)=\operatorname{Shape}(n')
 \iff
 \exists\rho\in\TRen(\Delta_n,\Delta_{n'}).\ K_n\Rightarrow_\rho K_{n'}.
\]
Either side puts $K_n$ and $K_{n'}$ in the same $\sim$-orbit, preserving
$\operatorname{out}(K_n)=\operatorname{out}(n)$. Because $Q(K_n)\sim K_n$, the
constructive selector yields a witness $\sigma_n\in\TRen(C_n,\Delta_n)$ and
composing with $\iota_n$ yields $\omega_n=\iota_n\circ\sigma_n$.\qed
\end{proof}

\paragraph{Why carrier completeness and effective support matter.}
If enumeration stops early under a heuristic budget instead of satisfying
carrier completeness, equal shapes still witness equivalence ($\Rightarrow$),
but two structurally equivalent expressions need not select the same minimum
($\Leftarrow$). For example, consider the open Alloy body \texttt{p in Ad} in contexts
$(u\!:\!\mathsf{User},p\!:\!\mathsf{Photo})$ and $(p\!:\!\mathsf{Photo})$.
In each context the body reads only $p$. Its effective support therefore
contains one photo slot, so the identity on that slot matches the two shapes.

\subsection{Building canonical records and search cost}
\label{sec:canonicalization-procedure}

A canonical record $(K_n,p_n,\sigma_n,\iota_n)$ stores the kernel $K_n$, least
shape $p_n=\operatorname{Shape}(n)$, witness $\sigma_n:C_n\cong\Delta_n$,
inclusion $\iota_n:\Delta_n\hookrightarrow\Gamma_n$, and extraction/normalization
provenance. Construction proceeds in four steps:
(1)~follow leaders, apply licensed local rules, and flatten visible matching
operators;
(2)~compute the effective support $\Delta_n$, forming $K_n$ and $\iota_n$;
(3)~for each typed bijection $\sigma\in\TRen(C_n,\Delta_n)$, apply
$\sigma^{-1}$ to $K_n$, normalize local invocation and binder orbits, and
aggregate bags or deduplicate sets;
(4)~select the least resulting shape $p_n$ and retain $\sigma_n$ and its
derivation.
Only $p_n$ acts as the collision key; the remaining fields place the shape in
$\Gamma_n$.

When a child merge or new symmetry triggers \emph{rebuilding}, affected parent
records are recomputed and colliding keys are checked against output types and
maps. Crucially, the exact enumeration in step~(3) searches
$\prod_\tau|\Delta_{n,\tau}|!$ typed bijections, where $\Delta_{n,\tau}$ is
the fibre of slots of type $\tau$. Slots of distinct types never mix in the
permutation search, whereas $N$ slots of a single type retain the full $N!$
search space—a scaling behavior we measure directly in Section~\ref{sec:eval-synthetic}.

\begin{figure}[t]
\centering
\scriptsize
\setlength{\fboxsep}{4pt}

\fbox{%
\begin{minipage}{0.94\linewidth}
\textbf{Model context}\par
\ttfamily
sig User \{ follows: set User, sees: set Photo, posts: set Photo \}\\
sig Photo \{ date: one Day \}\hspace{1.8em}%
sig Ad extends Photo \{\}\hspace{1.8em}sig Day \{\}
\end{minipage}%
}

\vspace{4pt}

\begin{minipage}[t]{0.47\linewidth}
\raggedright
\textbf{LEFT}\par
\ttfamily
pred left \{\\
\hspace*{1em}all u: User, p: Photo |\\
\hspace*{2em}p in u.sees implies\\
\hspace*{3em}(p not in Ad) implies\\
\hspace*{4em}p in u.follows.posts or p in Ad\\
\}
\end{minipage}
\hfill
\begin{minipage}[t]{0.47\linewidth}
\raggedright
\textbf{RIGHT}\par
\ttfamily
pred right \{\\
\hspace*{1em}all photo: Photo, user: User |\\
\hspace*{2em}photo in user.sees implies\\
\hspace*{3em}photo in Ad or\\
\hspace*{3em}photo in user.follows.posts\\
\}
\end{minipage}

\vspace{4pt}

\begin{minipage}{0.95\linewidth}
Let
\(A(x,y)\equiv y\in x.\mathsf{sees}\),
\(B(x,y)\equiv y\in x.\mathsf{follows}.\mathsf{posts}\), and
\(C(y)\equiv y\in\mathsf{Ad}\).
Under the canonical context
\(\Gamma_Q=\{q_0:\mathsf{User},\,q_1:\mathsf{Photo}\}\):
\[
\begin{aligned}
\mathsf{left}
&\stackrel{\mathrm{Imp},\,\alpha_{\Gamma_Q}}{\Longrightarrow}
  \neg A(q_0,q_1)\lor C(q_1)\lor B(q_0,q_1)\lor C(q_1)
\stackrel{\mathrm{ACI}_{\lor}}{\Longrightarrow} Q,
\\
\mathsf{right}
&\stackrel{\mathrm{Imp},\,\alpha_{\Gamma_Q},\,\pi_{\mathrm{bind}}}{\Longrightarrow}
  \neg A(q_0,q_1)\lor C(q_1)\lor B(q_0,q_1)
\stackrel{\mathrm{ACI}_{\lor}}{\Longrightarrow} Q.
\end{aligned}
\]
\end{minipage}

\vspace{2pt}

\fbox{%
\begin{minipage}{0.94\linewidth}
{\centering\textbf{Shared canonical form $Q$ (distance $0$ between \texttt{left} and \texttt{right})}\par}
\ttfamily
all \_q0: User, \_q1: Photo |\\
\hspace*{1em}\_q1 not in \_q0.sees or \_q1 in Ad or \_q1 in \_q0.follows.posts
\end{minipage}%
}

\vspace{3pt}

\begin{minipage}{0.95\linewidth}
\textbf{Inequivalent variant (repair distance $>0$):}\quad
{\ttfamily pred bug \{ all user: User, photo: Photo | photo in user.sees implies photo in user.follows.posts \}}
omits {\ttfamily photo in Ad}. Canonical equality fails ($d>0$), and repair
distance measures structural edits between the normalized projection of
\texttt{bug} and $Q$.
\end{minipage}

\caption{Running Alloy example. Implication elimination exposes flat
disjunctions, typed renaming and independent binder permutation align
\texttt{User} and \texttt{Photo} coordinates to $\Gamma_Q$, and
$\Set$-port normalization drops order and the duplicate \texttt{p in Ad}.}
\label{fig:alloy-certificate-example}
\end{figure}


Figure~\ref{fig:alloy-certificate-example} returns to our opening Alloy
example. Implication elimination exposes flat disjunctions, typed renaming and
independent binder permutation align the \texttt{User} and \texttt{Photo}
slots into $\Gamma_Q$, and $\Set$-port normalization removes operand order and
the duplicate \texttt{p in Ad}, yielding distance~0 between \texttt{left} and
\texttt{right} under the implemented rewrite theory. By contrast, \texttt{bug}
omits \texttt{photo in Ad}, so canonical equality fails ($d>0$), and repair
distance measures edits between its normalized repair projection and $Q$.

\section{An Algebraic Soundness Argument}
\label{sec:metatheory}

We now connect structural derivations, shape collisions, and graph updates to
the typed source equational theory $\mathcal T$, whose derivations
$\mathcal T\vdash_\Gamma s=t$ are generated by primitive laws, forward
congruence, reflexivity, symmetry, transitivity, and transport along embeddings
$e\in\TEmb(\Gamma,\Gamma')$. In every typed model $\mathfrak M$ validating the
primitive laws and congruence rules, environment restriction satisfies
$\operatorname{eval}_{\mathfrak M}(e\cdot t,\eta')
=\operatorname{eval}_{\mathfrak M}(t,\eta'\circ e)$, so every derivable
certificate holds semantically in $\mathfrak M$.

\subsection{From structural derivations to source equations}

A \emph{typed structural realizer} $R$ interprets port atoms, containers,
unary binders, binder blocks, and nodes as source terms and commutes with every
embedding $e$. Constructor laws supply one-step equational certificates for
$\Seq$ (positional bijection), $\Bag$ (occurrence bijection), $\Set$ (mutual
forward and backward matching maps), unary binders (scope extension),
descriptor-admissible binder-block permutations, and outer nodes, while leaf
steps carry individual equational certificates.

\begin{proposition}[Structural realization and congruence soundness]
\label{thm:structural-soundness}
Given a typed structural realizer $R$ commuting with embeddings, licensed
constructor congruence laws, and per-leaf equational certificates, every direct
structural derivation from $v$ to $v'$ under $e\in\TEmb(\Gamma,\Gamma')$ yields
\[
 \mathcal T\vdash_{\Gamma'} e\cdot R(v)=R(v'),
 \qquad
 \operatorname{eval}_{\mathfrak M}(e\cdot R(v),\eta)
 =\operatorname{eval}_{\mathfrak M}(R(v'),\eta)
\]
in every typed model $\mathfrak M$ and environment $\eta$. The same holds for
the indexed reflexive--symmetric--transitive closure under typed renamings
$\rho\in\TRen(\Gamma,\Gamma')$, including descriptor-admissible binder-block
permutations.
\end{proposition}
\begin{proof}
Structural induction gives naturality $R(e\cdot v)=e\cdot R(v)$ from the
supplied constructor commutation equations. Induction on the direct structural
derivation then constructs the equational certificate: leaf steps use their
supplied certificates; $\Seq$, $\Bag$, $\Set$, unary-binder, and binder-block
cases apply the corresponding licensed constructor law to the inductive child
certificates; and node steps apply node congruence.

For the closure indexed by renamings, reflexivity supplies the identity
certificate. Symmetry transports a $\rho$-indexed certificate along $\rho^{-1}$
and applies certificate symmetry; transitivity transports the first certificate
along $\rho_2$ and composes with the second, using naturality to match
$\rho_2\circ\rho_1$. Induction over the resulting finite certificate establishes
evaluation equality in every typed model $\mathfrak M$.\qed
\end{proof}

\paragraph{Why descriptor admissibility and constructor laws matter.}
A type-preserving permutation cannot be applied blindly to every binder or
operator. In Alloy, swapping \texttt{u: User, p: Photo} is admissible in
\texttt{all u: User, p: Photo | ...}, but violates meaning in the dependent
block \texttt{all u: User, p: u.sees | ...} or inside a set comprehension
\texttt{\{u: User, p: Photo | ...\}} whose column order is observable.
Likewise, under signed 4-bit Alloy arithmetic with overflow rejection,
$(7+1)+(-1)$ fails at the first sum while $7+(1+(-1))$ evaluates to~$7$, so
flattening requires an explicit operator-profile associativity law.

\subsection{From a shape collision to an equation}

A \emph{certified canonical record} over a shared canonical context $C$, output
type $\tau$, shape carrier, and realization function $\operatorname{real}$
comprises: (i)~a certified extraction $(n_i,K_i,\iota_i)$ with
$\iota_i\in\TEmb(\Delta_i,\Gamma_i)$ and finite replayable provenance from
$n_i$ to $\iota_i\cdot K_i$; (ii)~a typed support renaming
$\sigma_i\in\TRen(C,\Delta_i)$; and (iii)~finite replayable provenance from
$\sigma_i^{-1}\cdot K_i$ to $\operatorname{real}(p_i)$.

\begin{proposition}[Compatible effective-shape collisions preserve equality]
\label{thm:collision-soundness}
Suppose two certified canonical records have $p_1=p_2$, and set
$\rho=\sigma_2\circ\sigma_1^{-1}\in\TRen(\Delta_1,\Delta_2)$. Then
\[
 \mathcal T\vdash_{\Delta_2} \rho\cdot K_1 = K_2.
\]
Moreover, for any common context $\Omega$ and embeddings
$e_i\in\TEmb(\Gamma_i,\Omega)$ satisfying the compatibility condition
\begin{equation}
 e_1\circ\iota_1 = e_2\circ\iota_2\circ\rho,
 \label{eq:collision-compatible}
\end{equation}
we have $\mathcal T\vdash_\Omega e_1\cdot n_1 = e_2\cdot n_2$.
\end{proposition}
\begin{proof}
Replaying each finite provenance by induction (using primitive laws,
Proposition~\ref{thm:structural-soundness}, symmetry, and transitivity) yields
$\mathcal T\vdash_{\Gamma_i} n_i=\iota_i\cdot K_i$ and
$\mathcal T\vdash_C \sigma_i^{-1}\cdot K_i=\operatorname{real}(p_i)$. Since
$p_1=p_2$, $\operatorname{real}(p_1)=\operatorname{real}(p_2)$, so composing
the first normalization with the symmetric second gives
$\mathcal T\vdash_C \sigma_1^{-1}\cdot K_1=\sigma_2^{-1}\cdot K_2$.
Transporting along $\sigma_2$ yields $\mathcal T\vdash_{\Delta_2}\rho\cdot K_1=K_2$.

For the ambient equation, transport the extraction certificates along $e_1$ and
$e_2$, transport the kernel equation along $e_2\circ\iota_2$, and rewrite with
the compatibility condition~\eqref{eq:collision-compatible}:
\[
 e_1\cdot n_1
 =(e_1\circ\iota_1)\cdot K_1
 =(e_2\circ\iota_2\circ\rho)\cdot K_1
 =(e_2\circ\iota_2)\cdot K_2
 =e_2\cdot n_2,
\]
producing a valid certificate in $\Omega$.\qed
\end{proof}

\paragraph{Why ambient compatibility matters.}
Consider $\Omega=\{u,v:\mathsf{User},\,p:\mathsf{Photo}\}$ containing
$n_1=(p\in u.\mathsf{sees})$ and $n_2=(p\in v.\mathsf{sees})$. Both share the
canonical shape $q_1\in q_0.\mathsf{sees}$, inducing $\rho(u)=v$ and
$\rho(p)=p$. Under $e_1=e_2=\operatorname{id}_\Omega$, however,
$e_1(\iota_1(u))=u\neq v=e_2(\iota_2(\rho(u)))$, violating
\eqref{eq:collision-compatible}. The compatibility premise allows both
expressions to share a shape bucket without falsely equating $p\in u.\mathsf{sees}$
with $p\in v.\mathsf{sees}$.

\subsection{The graph invariant and finite-unfolding soundness}

Assign each class $a$ a witness term $w_a$ over $S_a$. A stored case $c$ of
$a$ has realization $r_c$ over $T_c$ and inclusion
$\iota_c\in\TEmb(S_a,T_c)$. A graph state is \emph{coherent} when every live
obligation carries a finite equational certificate in three families:
\begin{align}
 \text{(EC)}\qquad
 &\mathcal T\vdash_{T_c} r_c=\iota_c\cdot w_a
       &&\text{for each stored case }c\text{ of }a,\label{eq:EC}\\
 \text{(PC)}\qquad
 &\mathcal T\vdash_{S_a} w_a=p\cdot w_b
       &&\text{for each parent edge }a\xrightarrow{p}b,\label{eq:PC}\\
 \text{(SC)}\qquad
 &\mathcal T\vdash_{S_a} w_a=\pi\cdot w_a
       &&\text{for each symmetry generator }\pi\text{ of }a.\label{eq:SC}
\end{align}
We admit traces built from four evidence-preserving obligation transitions:
addition of an obligation with a local certificate; rekeying an existing
obligation via typed left and right endpoint bridges; composition of two
existing obligations via three endpoint bridges; and removal or inheritance of
already justified obligations. A \emph{certified finite unfolding} of $m*a$ to
$t$ in $\Gamma$ picks a stored case $c$ of $a$, an extension
$\bar m\in\TEmb(T_c,\Gamma)$ with $\bar m\circ\iota_c=m$, and finite provenance
relating $t$ to $\bar m\cdot r_c$. Focusing on finite unfoldings lets the
soundness theorem cover cyclic classes (such as $a=\{x,\,f(a)\}$) without
requiring coinductive semantics for infinite trees.

\begin{theorem}[Certificate-preserving reachability and finite-unfolding soundness]
\label{thm:semantic-sound}
Any finite trace of the four evidence-preserving obligation transitions
starting from a state with no live obligations ends in a coherent state.
Moreover, if two certified finite unfoldings in a caller context $\Gamma$
produce $t$ and $u$ from classes $a$ and $b$ with caller maps $m,m'$, if parent
paths reach a common leader $\ell$ with composite maps $\widehat m=m\circ p_a$
and $\widehat m'=m'\circ p_b$, and if there exists a finite word $\pi$ in the
recorded symmetry generators of $\ell$ (and their inverses) such that
$\widehat m=\widehat m'\circ\pi$, then:
\[
 \mathcal T\vdash_\Gamma t=u,
 \qquad
 \operatorname{eval}_{\mathfrak M}(t,\eta)
 =\operatorname{eval}_{\mathfrak M}(u,\eta)
\]
for every typed model $\mathfrak M$ and environment $\eta$, and for every
common context $\Omega$ and equal embeddings $e_1=e_2\in\TEmb(\Gamma,\Omega)$,
$\mathcal T\vdash_\Omega e_1\cdot t=e_2\cdot u$.
\end{theorem}
\begin{proof}
The empty state is vacuously coherent. By induction on the trace, addition
either inherits an existing certificate or uses the supplied local one;
rekeying transports the source certificate and composes it with the left and
right endpoint bridges; composition transports and chains two source
certificates across the three endpoint bridges; and removal or inheritance
retains existing certificates. Hence the final state satisfies (EC), (PC), and
(SC).

By induction on the parent path from $a$ to $\ell$, transporting each (PC)
certificate into $\Gamma$ gives $\mathcal T\vdash_\Gamma m\cdot w_a=\widehat m\cdot w_\ell$.
Replaying the unfolding provenance and transporting (EC) along $\bar m$ gives
$\mathcal T\vdash_\Gamma t=\bar m\cdot r_c=m\cdot w_a=\widehat m\cdot w_\ell$,
and likewise $\mathcal T\vdash_\Gamma u=\widehat m'\cdot w_\ell$. Induction on
the symmetry word $\pi$ using (SC), symmetry, transitivity, and transport
produces a certificate for $w_\ell=\pi\cdot w_\ell$. Transporting along
$\widehat m'$ and using $\widehat m=\widehat m'\circ\pi$ gives
\[
 \mathcal T\vdash_\Gamma t
 =\widehat m\cdot w_\ell
 =(\widehat m'\circ\pi)\cdot w_\ell
 =\widehat m'\cdot(\pi\cdot w_\ell)
 =\widehat m'\cdot w_\ell
 =u.
\]
Semantic equality in every typed model follows from certificate soundness, and
transporting along $e_1=e_2$ yields the common-context certificate in $\Omega$.\qed
\end{proof}

\paragraph{Why recorded (SC) symmetry certificates matter.}
Suppose leader $\ell$ represents $x\in y.\mathsf{follows}$ over
$\{x,y:\mathsf{User}\}$. In $\Gamma=\{a,b:\mathsf{User}\}$, the two unfoldings
$a\in b.\mathsf{follows}$ and $b\in a.\mathsf{follows}$ both reach $\ell$ and
their embeddings differ by a type-preserving slot swap. Because \texttt{follows}
is not symmetric, that swap has no (SC) certificate, so
Theorem~\ref{thm:semantic-sound} does not equate them; once an expression is
genuinely symmetric (such as $x\in y.\mathsf{follows}\land y\in x.\mathsf{follows}$),
an (SC) certificate justifies the swap and the theorem validates the merge.
